% ============================================================
\section{Discussion}
\label{sec:discussion}
% ============================================================

\subsection{Improvements Over km/kmtools: A Systematic Assessment}

\kmerdet{} addresses each of the ten limitations identified in the thesis
validation study~\cite{kamtool}.  This subsection evaluates each improvement
in the context of the original findings.

\subsubsection{SNV Sensitivity (Baseline: 77\%)}

The thesis found that 23\% of true SNVs were missed, with false negatives
concentrated at VAF $< 0.1\%$ and in high-background-noise genomic regions.
Three complementary improvements in \kmerdet{} address this:

\textbf{Adaptive thresholds} (Section~\ref{sec:adaptive}): Rather than the
fixed \texttt{count=2} floor, the Poisson noise model sets the threshold
proportional to local coverage.  At $50,000\times$ coverage (achievable with
modern ultra-deep targeted panels), a variant at $0.1\%$ VAF has expected
count $\sim$50, and the error background has expected count $\sim$50 as well
(at $0.1\%$ per-base error rate).  A fixed threshold of 2 would pass all error
branches; an adaptive threshold of $\sim$20 (the $99.9\%$ quantile of
$\text{Poisson}(16.7)$) reduces false extensions while detecting variants above
$0.1\%$ VAF.

\textbf{Bidirectional walking}: Variants within $k$ positions of the target
start or end were systematically missed by forward-only walking.  Backward
walking from the target end provides symmetric coverage of boundary regions.

\textbf{Statistical quality scores}: The Phred-scaled QUAL score enables
downstream filtering by statistical significance rather than by fixed count
thresholds alone, allowing borderline calls (QUAL 20--29) to be reviewed
separately from high-confidence calls (QUAL $\geq 30$).

\subsubsection{INDEL Sensitivity (Baseline: 38\% Overall, $\sim$3\% for $>7$~bp)}

The INDEL sensitivity limitation is fundamentally rooted in the k-mer length
constraint: for $k = 31$, insertions of length $L > 31$ have zero spanning
k-mers.  \kmerdet{}'s multi-k strategy directly addresses this.

For insertions of length $L$, the number of spanning k-mers with k-mer length
$k$ is $\max(0, k - L)$.  The number of spanning k-mers \emph{triples} when
switching from $k = 31$ to $k = 21$ for a 20~bp insertion: from
$\max(0, 31 - 20) = 11$ to $\max(0, 21 - 20) = 1$ --- still small but
non-zero.  For an 18~bp insertion, $k = 21$ provides $21 - 18 = 3$ spanning
k-mers versus $31 - 18 = 13$ at $k = 31$.  The absolute count of these spanning
k-mers at 0.1\% VAF and 3000$\times$ coverage is $\sim$3, which is above the
minimum count floor of 2.

The multi-k consensus voting approach (Section~\ref{sec:multik}) combines evidence
from $k \in \{21, 31, 41\}$, providing complementary sensitivity profiles:
$k = 21$ is most sensitive for large indels; $k = 41$ provides highest
specificity in repetitive regions; $k = 31$ is the validated standard for SNVs.

\subsubsection{VAF Detection Threshold (Baseline: $\sim$0.1\%)}

Pushing the detection threshold below $0.1\%$ VAF requires either higher
sequencing depth, better error suppression, or smarter statistical modelling.
\kmerdet{} addresses the last of these through the negative binomial model and
bootstrap confidence intervals.  The PoN filter reduces the false positive
background, allowing lower effective detection thresholds without inflating the
false positive rate.

The fundamental barrier remains sequencing depth and error rate: at $0.01\%$
VAF with $5000\times$ coverage, variant k-mers have expected count $\sim$0.5,
below any reasonable minimum count threshold.  Duplex consensus sequencing,
which the \kam{} pipeline currently does not use for k-mer counting, could push
the error rate from $\sim$$10^{-4}$ to $\sim$$10^{-7}$, enabling reliable
detection at $0.01\%$ VAF.

\subsubsection{Adaptive Parameters}

The \km{} pipeline applies the same \texttt{count=2, ratio=0.00001} to all
targets regardless of local coverage.  \kmerdet{}'s adaptive mode computes
per-target thresholds from the observed reference k-mer distribution, effectively
implementing the per-sample parameter table recommended in the thesis (count=5
at $5000\times$, count=10 at $>10,000\times$) automatically and continuously
rather than through discrete manual tiers.

\subsubsection{Target Design Failures}

Non-unique anchor k-mers cause silent walking failures in \km{}: the tool
produces output, but the output may be based on k-mers from unrelated genomic
regions that happen to share the anchor sequence.  \kmerdet{} performs anchor
uniqueness validation before each target is processed, flagging non-unique
anchors with a warning and optionally extending the flank until a unique anchor
is found.

\subsubsection{Graph Construction Limitations}

The fixed edge weighting scheme (0.01 for reference, 1.0 for non-reference) does
not reflect the actual coverage difference between alleles.  A future improvement
is coverage-weighted edges, where the edge weight is inversely proportional to
the k-mer count: $w = 1 / \log(\text{count} + 1)$.  This would guide Dijkstra
toward high-coverage (high-confidence) paths.  \kmerdet{}'s graph module is
designed to support alternative weighting functions through a configurable
weight function, but the fixed weighting remains the default for backward
compatibility.

\subsubsection{Quantification Model}

The NNLS model assumes homoscedastic errors (equal variance for all k-mer counts),
which does not hold in practice: high-count k-mers have higher absolute variance
(following Poisson scaling).  Weighted NNLS, with weights $w_i = 1/\hat{c}_i$,
would provide more accurate rVAF estimates by downweighting the most variable
observations.  \kmerdet{}'s bootstrap CI computation implicitly captures some
of this heteroscedasticity through Poisson resampling, but the point estimate
remains unweighted for backward compatibility.

\subsection{Remaining Limitations}
\label{sec:discussion-limitations}

Despite the improvements described above, several fundamental limitations remain:

\textbf{Targeted-only analysis}: \kmerdet{} requires pre-designed target
sequences and cannot perform genome-wide de novo variant discovery.  This is
appropriate for the tumor-informed MRD monitoring use case but limits the tool's
applicability to screening contexts where the set of variants is not known in advance.

\textbf{INDEL representation instability}: Even with INDEL normalisation, the
k-mer walking approach can produce different representations of the same INDEL
depending on the target sequence design and the local sequence context.  Robust
equivalence testing across different target designs remains an open problem.

\textbf{No strand-bias annotation}: \jf{}'s canonical counting mode merges
forward and reverse strand k-mers.  Strand bias is a powerful filter for
Illumina sequencing artefacts (single-strand artefacts produce k-mers from
only one strand).  Implementing strand-aware k-mer counting would require either
two separate jellyfish databases (one per strand) or a custom k-mer counter,
and is the single most impactful quality annotation missing from the current
implementation.

\textbf{Large deletions}: Deletions of $>k$~bp produce no junction k-mers that
span the breakpoint; only the absence of reference k-mers signals the deletion.
While the graph pathfinding can theoretically detect this absence (the reference
path has gaps), the current implementation does not explicitly report large
deletions.  The 2-kupl approach~\cite{pelletier2021kupl}, which detects variants
by comparing k-mer populations between matched samples, is better suited to large
deletion detection.

\textbf{No germline distinction}: Without a matched normal sample (buffy coat),
\kmerdet{} cannot distinguish germline heterozygous variants from somatic
variants at similar VAF.  The PoN filter partially addresses this for common
germline variants, but rare germline variants will be reported as somatic findings
unless a matched normal analysis is available.

\subsection{Clinical Implications}
\label{sec:discussion-clinical}

The improvements in \kmerdet{} strengthen the case for k-mer-based variant
detection as a rapid preliminary screening tool in liquid biopsy clinical
workflows.  The original \kam{} pipeline demonstrated that k-mer analysis
could provide same-day preliminary results ($\sim$6 minutes) before the
standard alignment-based pipeline completes, with 77\% SNV sensitivity.

With \kmerdet{}'s projected improvements, the sensitivity gap relative to
alignment-based callers (77\% vs.\ $\sim$90\%) narrows substantially.  For MRD
monitoring applications where the variant panel is fixed and known from the
primary tumour biopsy, the remaining sensitivity gap may be acceptable given
the dramatic speed advantage.  A hybrid strategy --- k-mer preliminary
screening followed by alignment-based confirmation of positive calls --- remains
the clinically recommended approach:

\begin{enumerate}
  \item \textbf{Primary screen} (\kmerdet{}, $\sim$4~min): Report all variants
    with QUAL $\geq 20$ as preliminary.
  \item \textbf{Confirmation} (BWA + Mutect2, $\sim$1~h): Confirm positive calls
    and resolve equivocal results.
  \item \textbf{Longitudinal tracking} (\kmerdet{}, $<5$~min per sample):
    Once a variant is confirmed, serial monitoring with \kmerdet{} provides
    rapid rVAF trends.
\end{enumerate}

The bootstrap confidence intervals on rVAF are particularly valuable for
longitudinal tracking: a 10-sample patient time series can display rVAF trends
with uncertainty bounds, enabling early detection of the rising trend that
precedes clinical relapse.

The clinical utility of CHIP suppression via the PoN filter cannot be
overstated: recent large-scale data suggest that $>50\%$ of cancer patients
have at least one CHIP mutation in a reportable clinical gene~\cite{blasco2015chip},
making CHIP filtering a prerequisite for clinical-grade reporting.

\subsection{Future Work}
\label{sec:discussion-future}

Several high-impact improvements are planned for future \kmerdet{} versions:

\subsubsection{UMI-Aware K-mer Counting}

Current k-mer counting is at the read level: multiple PCR copies of the same
original molecule contribute equally to k-mer counts.  A molecule-level count
(number of distinct UMI families supporting each k-mer) would be intrinsically
immune to PCR amplification bias.  Implementation requires either a custom
k-mer counter that tracks UMI-to-k-mer associations or a streaming integration
with the HUMID deduplication step.

\subsubsection{In-Process Deduplication}

HUMID currently runs as a separate process that writes intermediate FASTQ files
to disk.  Integrating UMI-aware deduplication into \kmerdet{}'s preprocessing
stage would eliminate this disk I/O, reducing the bottleneck from $\sim$2~min to
near zero and enabling a sub-2-minute end-to-end pipeline.

\subsubsection{Coverage-Weighted Pathfinding}

Replacing fixed edge weights (0.01/1.0) with count-derived weights
$w = 1/\log(\text{count} + 1)$ would guide Dijkstra toward high-confidence
variant paths, improving pathfinding accuracy in complex regions with multiple
competing paths.

\subsubsection{Maximum Likelihood rVAF Estimation}

The NNLS decomposition minimises a squared-error objective that implicitly
assumes Gaussian noise.  A Poisson or negative binomial likelihood model would
provide more statistically principled rVAF estimates and naturally incorporate
heteroscedasticity (higher variance for higher counts).

\subsubsection{Deep Learning Integration}

A gradient-boosted classifier trained on features extracted from \kmerdet{}
variant calls (rVAF, min\_coverage, GC content, homopolymer proximity,
linguistic complexity, path count, coverage uniformity) could substantially
reduce false positives without sacrificing true positive sensitivity.  The
\texttt{linfa} Rust crate provides gradient-boosted tree implementations that
can be trained offline and deployed as a $<1$~MB serialised model shipped with
the binary.

\subsubsection{Multi-Sample Joint Analysis}

Processing multiple samples from the same patient simultaneously (longitudinal
analysis) or from multiple patients sharing the same target panel (cohort
analysis) enables borrowing of statistical strength across samples.  A variant
consistently present in serial blood draws is more credible than a single-sample
observation; a variant present in only one patient from a cohort with no history
of the same variant is more likely to be a true somatic mutation than a
recurrent artifact.

\subsubsection{Native K-mer Counting}

Replacing the \jf{} dependency with a native Rust k-mer counter would eliminate
the last external dependency and enable UMI-aware counting as a built-in
feature.  The counting algorithm (lock-free hash table with CAS operations) is
well-understood and implementable in safe Rust using \texttt{dashmap} or a
custom \texttt{AtomicU64}-based structure.
